% ==========
% CHAPTER 2
% ==========
\OpenGrid

\ParallelChapter
  {Reasoning Habits}
  {Thói quen luận}
  {}

\TriPara
  {
    \EN{
      \begin{flushright}
      \emph{Reasoning and sense making should be a part of the mathematics classroom every day.}
      \end{flushright}}
  }
  {
    \VI{
      \begin{flushright}
      \emph{Luận và hiểu phải hiện diện trong lớp học toán mỗi ngày.}
      \end{flushright}}
  }
  {
    \NOTE{Xác lập lập trường trung tâm của chương: luận và hiểu không phải phần phụ trợ, mà phải hiện diện trong thực hành dạy học toán hằng ngày.}
  }

\TriPara
  {
    \EN{A focus on reasoning and sense making implies that ``covering'' mathematical topics is not enough. Students also need to experience and develop mathematical reasoning habits (cf. Cuoco, Goldenberg, and Mark [1996]; Driscoll [1999]; Pólya [1952, 1957]; Schoenfeld [1983]; Harel and Sowder [2005]). A \emph{reasoning habit} is a productive way of thinking that becomes common in the processes of mathematical inquiry and sense making. The following list of reasoning habits illustrates the types of thinking that should become routine and fully expected in the classroom culture of all mathematics classes across all levels of high school.}
  }
  {
    \VI{Việc lấy luận và hiểu làm trọng tâm hàm ý rằng chỉ ``dạy cho hết'' các chủ đề toán học là không đủ. Học sinh còn cần được trải nghiệm và phát triển các thói quen luận toán học (so sánh Cuoco, Goldenberg, and Mark [1996]; Driscoll [1999]; Pólya [1952, 1957]; Schoenfeld [1983]; Harel and Sowder [2005]). Một \emph{thói quen luận} là cách tư duy hiệu quả, dần trở nên quen thuộc trong quá trình tìm tòi toán học và hiểu. Danh sách thói quen luận dưới đây minh họa những kiểu tư duy cần trở thành thông lệ và được xem là điều đương nhiên trong văn hóa lớp học của mọi lớp toán ở mọi cấp độ trung học phổ thông.}
  }
  {
    \NOTE{Nếu lấy luận và hiểu làm trọng tâm, học sinh phải thường xuyên được tham gia vào những cách tư duy có hiệu quả trong học toán. Vì vậy, thói quen luận không phải phần phụ thêm sau nội dung, mà là những cách nghĩ cần được nuôi dưỡng đến mức trở thành thông lệ trong văn hóa lớp học toán ở bậc trung học phổ thông.}
  }

\TriPara
  {
    \EN{Approaching the list as a new set of topics to be taught in an already crowded curriculum is not likely to have the desired effect. Instead, attention to reasoning habits needs to be integrated within the curriculum to ensure that students both understand and can use what they are taught.}
  }
  {
    \VI{Việc xem danh sách thói quen luận như một tập hợp chủ đề mới cần đưa vào giảng dạy trong chương trình vốn đã quá tải khó có thể đem lại hiệu quả mong muốn. Thay vào đó, việc chú ý đến thói quen luận cần được tích hợp ngay trong chương trình để bảo đảm rằng học sinh vừa hiểu vừa có thể vận dụng những gì mình được học.}
  }
  {
    \NOTE{Ý ở đây là thói quen luận không nên được thêm vào như một mảng nội dung riêng trong chương trình đã chật, mà cần được nuôi dưỡng ngay trong cách tổ chức và dạy nội dung hiện có. Trọng tâm không phải là thêm chủ đề mới, mà là làm cho học sinh hiểu và dùng được những gì mình học.}
  }

\TriPara
  {
  \EN{
  \begin{itemize}
    \item  \emph{Analyzing a problem}, for example,
    \begin{itemize}
      \item \emph{identifying relevant} mathematical \emph{concepts, procedures, or representations} that reveal important information about the problem and contribute to its solution (for example, choosing a model for simulating a random experiment);
      
      \item \emph{defining relevant variables and conditions} carefully, including units if appropriate;
      
      \item \emph{seeking patterns and relationships} (for example, systematically examining cases or creating displays for data);
      
      \item \emph{looking for hidden structure} (for example, drawing auxiliary lines in geometric figures or finding equivalent forms of expressions that reveal different aspects of a problem);
      
      \item \emph{considering special cases or simpler analogs};
      
      \item \emph{applying previously learned concepts} to new problem situations, adapting and extending as necessary;
      
      \item \emph{making preliminary deductions and conjectures}, including predicting what a solution to a problem might involve or putting constraints on solutions; and
      
      \item \emph{deciding whether a statistical approach is appropriate}.
    \end{itemize}

    \item \emph{Implementing a strategy}, for example,
    \begin{itemize}
      \item \emph{making purposeful use of procedures};
      
      \item \emph{organizing the solution}, including calculations, algebraic manipulations, and data displays;
      
      \item \emph{making logical deductions} based on current progress, verifying conjectures, and extending initial findings; and
      
      \item \emph{monitoring progress toward a solution}, including reviewing a chosen strategy and other possible strategies generated by oneself or others.
    \end{itemize}

    \item \emph{Seeking and using connections} across different mathematical domains, different contexts, and different representations.
    
    \item \emph{Reflecting on a solution} to a problem, for example,
    \begin{itemize}
      \item \emph{interpreting a solution} and how it answers the problem, including making decisions under uncertain conditions;
      
      \item \emph{considering the reasonableness of a solution}, including whether any numbers are reported at an unreasonable level of accuracy;
      
      \item \emph{revisiting initial assumptions} about the nature of the solution, including being mindful of special cases and extraneous solutions;
      
      \item \emph{justifying or validating a solution}, including through proof or inferential reasoning;\vspace{\baselineskip}
      
      \item \emph{recognizing the scope of inference} for a statistical solution;
      
      \item \emph{reconciling different approaches} to solving the problem, including those proposed by others;
      
      \item \emph{refining arguments} so that they can be effectively communicated; and
      
      \item \emph{generalizing a solution} to a broader class of problems and looking for connections with other problems.
    \end{itemize}
  \end{itemize}
  }
  }
  {
  \VI{
  \begin{itemize}
    \item \emph{Phân tích bài toán}, chẳng hạn:
    \begin{itemize}
      \item \emph{xác định khái niệm, thủ tục hoặc biểu diễn toán học thích đáng}---những yếu tố làm lộ thông tin quan trọng về bài toán và góp phần vào lời giải (ví dụ, chọn mô hình để mô phỏng thí nghiệm ngẫu nhiên);\vspace*{\baselineskip}

      \item \emph{xác định cẩn thận biến và điều kiện thích đáng}, bao gồm đơn vị đo nếu phù hợp;

      \item \emph{tìm kiếm mẫu hình và quan hệ} (ví dụ, xét trường hợp một cách có hệ thống hoặc tạo biểu đồ cho dữ liệu);

      \item \emph{tìm kiếm cấu trúc ẩn} (ví dụ, kẻ đường phụ trong hình học hoặc tìm dạng tương đương của biểu thức để làm lộ những khía cạnh khác nhau của bài toán);

      \item \emph{xét trường hợp đặc biệt hoặc dạng tương tự đơn giản hơn};

      \item \emph{vận dụng khái niệm đã học} vào tình huống bài toán mới, điều chỉnh và mở rộng khi cần;\vspace{\baselineskip}

      \item \emph{đưa ra suy diễn và phỏng đoán sơ bộ}, bao gồm dự đoán lời giải có thể liên quan đến điều gì hoặc đặt ràng buộc cho lời giải; và\vspace{\baselineskip}

      
      \item \emph{quyết định liệu cách tiếp cận thống kê có phù hợp hay không}.
    \end{itemize}

    \item \emph{Triển khai chiến lược}, chẳng hạn:
    \begin{itemize}
      \item \emph{vận dụng thủ tục có chủ đích};

      \item \emph{tổ chức lời giải}, bao gồm phép tính, biến đổi đại số và biểu diễn dữ liệu;\vspace{\baselineskip}

      \item \emph{thực hiện suy diễn logic} dựa trên tiến triển hiện tại, kiểm tra phỏng đoán và mở rộng kết quả ban đầu; và

      \item \emph{theo dõi tiến trình đi đến lời giải}, bao gồm xem xét lại chiến lược đã chọn và chiến lược khả dĩ khác do chính mình hoặc người khác đề xuất.
    \end{itemize}

    \item \emph{Tìm kiếm và sử dụng kết nối} giữa những lĩnh vực toán học khác nhau, bối cảnh khác nhau và biểu diễn khác nhau.

    \item \emph{Phản tư về lời giải}, chẳng hạn:\vspace{\baselineskip}

    \begin{itemize}
      \item \emph{diễn giải lời giải} và cách lời giải trả lời bài toán, bao gồm việc đưa ra quyết định trong điều kiện không chắc chắn;

      \item \emph{xem xét tính hợp lí của lời giải}, bao gồm việc đánh giá liệu số liệu được báo cáo có ở mức chính xác bất hợp lí hay không;\vspace{\baselineskip}

      \item \emph{xem lại giả định ban đầu} về bản chất của lời giải, bao gồm việc chú ý đến trường hợp đặc biệt và nghiệm ngoại lai;

      \item \emph{biện minh hoặc kiểm chứng lời giải}, bao gồm thông qua chứng minh hoặc suy luận thống kê;

      \item \emph{nhận diện phạm vi suy luận} đối với lời giải thống kê; 

      \item \emph{dung hòa những cách tiếp cận khác nhau} để giải bài toán, bao gồm cả cách do người khác đề xuất;

      \item \emph{tinh chỉnh lập luận} để có thể giao tiếp hiệu quả; và
      
      \item \emph{khái quát hóa lời giải} sang lớp bài toán rộng hơn và tìm kiếm kết nối với bài toán khác.
    \end{itemize}
  \end{itemize}
  }
  }
  {
    \NOTE{Thói quen luận không phải là kĩ thuật rời rạc, mà là nếp hoạt động trí tuệ cần hiện diện xuyên suốt quá trình học: từ phân tích bài toán, triển khai chiến lược, tìm kiếm và sử dụng kết nối, đến phản tư về lời giải. Trọng tâm không nằm ở việc làm cho xong bài toán, mà ở việc học sinh dần biết làm rõ bản chất vấn đề, điều khiển tiến trình giải, liên kết ý tưởng toán học trong bối cảnh và biểu diễn khác nhau, rồi xem xét lại lời giải để đánh giá, kiểm tra và mở rộng khi cần. Qua đó, luận được đặt như chuẩn thường nhật của văn hóa lớp học toán, chứ không phải yêu cầu phụ được thêm vào sau khi đã ``dạy hết'' nội dung được xem là chính.}
  }

\TriPara
  {
    \EN{Many of these reasoning habits could be construed to fit in more than one category, and students are expected to move naturally among various reasoning habits as they are needed. Section 2 of this publication offers examples of how reasoning habits can be promoted in the high school classroom, with specific references to the reasoning habits listed above.}
  }
  {
    \VI{Nhiều thói quen luận trong số này có thể được xem là thuộc hơn một nhóm, và học sinh được kì vọng sẽ chuyển đổi tự nhiên giữa nhiều thói quen luận khi cần. Phần 2 của ấn phẩm này đưa ra ví dụ về cách thúc đẩy thói quen luận trong lớp học toán trung học phổ thông, với tham chiếu cụ thể đến những thói quen luận đã liệt kê ở trên.}
  }{\NOTE{}}

\ParallelSection
  {Progression of Reasoning}
  {Tiến trình phát triển của luận}
  {}

\TriPara
  {
    \EN{When reasoning is interwoven with sense making, and when teachers provide the necessary support and formative feedback, students can be expected to demonstrate growing levels of formality in their reasoning in the classroom, in their oral and written work, and in assessments throughout the high school years. Reasoning and sense making in the high school mathematics classroom require increasing levels of understanding, as outlined in the following:

    \begin{description}
      \item[empirical] --- the role of empirical evidence that supports, but does not justify, a conjecture---``it works in a number of cases'';
      \item[preformal] --- the role of intuitive explanations and partial arguments that lend insight into what is happening; and
      \item[formal] --- the role of formal argumentation (based on logic) in determining mathematical certainty (proof) or in making statistical inferences.
    \end{description}
    }
  }
  {
    \VI{Khi luận được đan xen với hiểu, và khi giáo viên cung cấp hỗ trợ cần thiết cùng phản hồi quá trình, học sinh có thể được kì vọng sẽ thể hiện mức độ hình thức ngày càng cao trong luận của mình---trong lớp học, trong sản phẩm nói và viết, và trong các hình thức đánh giá---xuyên suốt những năm trung học phổ thông. Luận và hiểu trong lớp học toán trung học phổ thông đòi hỏi mức độ hiểu ngày càng cao, như được phác thảo dưới đây:

    \begin{description}
      \item[thực nghiệm] --- vai trò của bằng chứng thực nghiệm là hỗ trợ, nhưng không biện minh cho phỏng đoán---``nó đúng trong một số trường hợp'';
      \item[tiền hình thức] --- vai trò của giải thích trực giác và lập luận bộ phận là soi sáng điều đang diễn ra; và
      \item[hình thức] --- vai trò của lập luận hình thức (dựa trên logic) trong việc xác lập tính chắc chắn toán học (chứng minh) hoặc thực hiện suy luận thống kê.
    \end{description}
    }
  }{\NOTE{}}

\TriPara
  {
    \EN{Thus, these levels show progress from less formal reasoning to more formal reasoning. However, each level has value. Students may continually shift among these levels, even within the same mathematical context. This shifting among levels is not only expected but desirable as students make sense of the context and reason their way to a conclusion. However, teachers play an essential role in encouraging students to explore more sophisticated levels of reasoning and sense making.}
  }
  {
    \VI{Do đó, các mức độ này cho thấy sự tiến triển từ luận ít hình thức hơn đến luận hình thức hơn. Tuy vậy, mỗi mức độ đều có giá trị riêng. Học sinh có thể liên tục chuyển đổi giữa những mức độ này, thậm chí trong cùng một bối cảnh toán học. Sự chuyển đổi giữa các mức độ không chỉ được kì vọng mà còn đáng mong đợi khi học sinh hiểu bối cảnh và luận để đi đến kết luận. Tuy nhiên, giáo viên đóng vai trò thiết yếu trong việc khuyến khích học sinh khám phá những mức độ tinh vi hơn của luận và hiểu.}
  }{\NOTE{}}

\TriPara
  {
    \EN{Formal argumentation includes both the ability of students to create meaningful chains of logical reasoning based on certain assumptions, definitions, and prior results, as well as the ability to read and evaluate reasoning given by others. The ability to determine the validity of an argument is important, as is an understanding of what the argument says about the ideas under consideration.}
  }
  {
    \VI{Lập luận hình thức bao gồm cả khả năng học sinh tạo ra chuỗi luận logic có ý nghĩa dựa trên giả định, định nghĩa và kết quả đã có, lẫn khả năng đọc và đánh giá luận do người khác đưa ra. Khả năng xác định tính hợp lệ của lập luận là quan trọng, cũng như việc hiểu lập luận ấy nói gì về ý tưởng đang xét.}
  }{\NOTE{}}

\ParallelSection
  {Developing Reasoning Habits in the Classroom}
  {Phát triển các thói quen luận trong lớp học}
  {}

\TriPara
  {
    \EN{Teachers can help students progress to higher levels of reasoning through judicious selection of tasks and the use of probing questions. Students can then learn to analyze their approach to solving problems, recognize the strengths and shortcomings of their current approach, and use the power of more formal reasoning to better formulate and justify mathematical conclusions. The continuing development of mathematical reasoning habits should be a priority in the high school classroom. The following is a preliminary list of tips for developing these habits.}
  }
  {
    \VI{Giáo viên có thể giúp học sinh tiến tới mức luận cao hơn thông qua việc lựa chọn nhiệm vụ thận trọng và sử dụng câu hỏi gợi mở. Khi đó, học sinh có thể học cách phân tích cách tiếp cận của mình khi giải bài toán, nhận ra điểm mạnh và hạn chế trong cách tiếp cận hiện tại, đồng thời sử dụng sức mạnh của luận hình thức hơn để hình thành và biện minh tốt hơn cho kết luận toán học. Việc liên tục phát triển thói quen luận toán học cần được xem là ưu tiên trong lớp học toán trung học phổ thông. Dưới đây là danh sách ban đầu gồm các gợi ý nhằm phát triển những thói quen này.}
  }{\NOTE{}}

\TriPara
  {
    \EN{
      \begin{itemize}
        \item Provide tasks that require students to figure things out for themselves.
        
        \item Ask students to restate the problem in their own words, including any assumptions they have made.
        
        \item Give students time to analyze a problem intuitively, explore the problem further by using models, and then proceed to a more formal approach.
        
        \item Resist the urge to tell students how to solve a problem when they become frustrated; find other ways to support students as they think and work.
        
        \item Ask students questions that will prompt their thinking---for example, ``Why does this work?'' or ``How do you know?''
        
        \item Provide adequate wait time after a question for students to formulate their own reasoning.
        
        \item Encourage students to ask probing questions of themselves and one another.
        
        \item Expect students to communicate their reasoning to their classmates and the teacher, orally and in writing, through using proper mathematical vocabulary.
        
        \item Highlight exemplary explanations, and have students reflect on what makes them effective.\vspace{\baselineskip}
        
        \item Establish a classroom climate in which students feel comfortable sharing their mathematical arguments and critiquing the arguments of others in a productive manner.
      \end{itemize}
    }
  }
  {
    \VI{
      \begin{itemize}
        \item Cung cấp nhiệm vụ đòi hỏi học sinh tự tìm ra cách tiếp cận.

        \item Yêu cầu học sinh diễn đạt lại bài toán bằng lời của mình, bao gồm cả giả định các em đã đưa ra.

        \item Cho học sinh thời gian phân tích bài toán bằng trực giác, khám phá sâu hơn thông qua mô hình, rồi sau đó mới chuyển sang cách tiếp cận hình thức hơn.

        \item Kiềm chế mong muốn chỉ ngay cho học sinh cách giải khi các em cảm thấy bế tắc; thay vào đó, hãy tìm cách hỗ trợ khác trong khi học sinh suy nghĩ và làm việc.

        \item Đặt câu hỏi gợi mở nhằm khơi gợi suy nghĩ của học sinh---chẳng hạn, ``Vì sao cách này đúng?'' hoặc ``Làm sao em biết?''.

        \item Dành đủ thời gian chờ sau mỗi câu hỏi để học sinh tự hình thành luận của mình.

        \item Khuyến khích học sinh đặt câu hỏi đào sâu cho chính mình và người khác.

        \item Kì vọng học sinh giao tiếp về luận của mình với bạn học và giáo viên, bằng lời nói và bằng văn viết, thông qua việc sử dụng đúng thuật ngữ toán học.

        \item Nêu bật những lời giải thích mẫu mực và yêu cầu học sinh phản tư về điều khiến chúng có hiệu quả.
        
        \item Xây dựng môi trường lớp học nơi học sinh cảm thấy thoải mái chia sẻ lập luận toán học của mình và phản biện lập luận của người khác theo cách mang tính xây dựng.
    \end{itemize}
    }
  }
  {
    \NOTE{Danh sách này quy tụ những thực hành lớp học có tác dụng nuôi dưỡng luận ở học sinh. Điểm chung là không đẩy học sinh thẳng tới lời giải, mà tạo điều kiện để các em tự phân tích tình huống, tự diễn đạt cách hiểu, tự nêu giả định, tự chất vấn, rồi từng bước chuyển từ trực giác và mô hình sang cách tiếp cận chặt chẽ hơn.

    Ở đây, luận không được xem như sản phẩm chỉ xuất hiện ở cuối cùng, mà như hoạt động được hình thành ngay trong quá trình học sinh nghĩ, nói, viết, hỏi và phản biện. Vì vậy, vai trò của giáo viên, trước hết, không phải là chỉ cách làm, mà là thiết kế nhiệm vụ, đặt câu hỏi, giữ nhịp chờ, và xây dựng môi trường lớp học đủ an toàn để lập luận có thể được đưa ra, được xem xét, và được cải thiện.}
  }

\TriPara
  {
    \EN{Teachers should refer to other resources for specific tasks, suggestions for questioning techniques, and so forth. NCTM and other professional organizations offer a number of valuable resources, both in print and online. Additional teaching ideas related to reasoning and sense making are also presented in more detail in the topic books that support this publication.}
  }
  {
    \VI{Giáo viên nên tham khảo các nguồn tư liệu khác để tìm nhiệm vụ cụ thể, gợi ý về kĩ thuật đặt câu hỏi, v.v. NCTM và nhiều tổ chức nghề nghiệp khác cung cấp nhiều nguồn tư liệu giá trị, cả bản in lẫn trực tuyến. Những ý tưởng dạy học bổ sung liên quan đến luận và hiểu cũng được trình bày chi tiết hơn trong các sách chuyên đề hỗ trợ ấn phẩm này.}
  }{\NOTE{}}

\ParallelSection
  {Reasoning as the Foundation of Mathematical Competence}
  {Luận như nền tảng của năng lực toán học}
  {}

\TriPara
{
\EN{Reasoning and sense making are inherent in developing mathematical competence, as discussed in \emph{Adding It Up} (Kilpatrick, Swafford, and Findell 2001); see figure 2.1. Sense making and conceptual understanding are closely interrelated. Procedural fluency includes learning with understanding and knowing which procedure to choose, when to choose it, and for what purpose. In the absence of reasoning, students may carry out procedures correctly but may also capriciously invoke incorrect or baseless rules, such as ``the square root of a sum is the sum of the square roots.'' They come to view procedures as steps they are told to do rather than a series of steps chosen for a specific purpose and based on mathematical principles. Without developing an understanding of procedures rooted in reasoning and sense making, students may be able to correctly perform those procedures but may think of them only as a list of ``tricks.'' As a result, they may have difficulty selecting an appropriate procedure to use in a given problem, or their seeming competence with simple tasks may evaporate in more complicated situations. Genuine procedural fluency requires both mastering technical skills and developing the understanding needed for using them appropriately. Strategic competence and adaptive reasoning are both directly addressed in the reasoning habits.}
}
{
  \VI{Luận và hiểu gắn liền với quá trình phát triển năng lực toán học, như đã được bàn trong \emph{Adding It Up} (Kilpatrick, Swafford, and Findell 2001); xem Hình 2.1. Hiểu và hiểu khái niệm có quan hệ chặt chẽ với nhau. Thành thạo thủ tục bao gồm việc học đi đôi với hiểu, cũng như biết nên chọn thủ tục nào, khi nào chọn thủ tục ấy, và chọn để làm gì. Khi thiếu luận, học sinh có thể thực hiện đúng thủ tục, nhưng cũng có thể tùy tiện viện đến quy tắc sai hoặc vô căn cứ, chẳng hạn như ``căn bậc hai của một tổng bằng tổng các căn bậc hai''. Các em dần xem thủ tục như những bước được bảo phải làm, thay vì như một chuỗi bước được lựa chọn cho mục đích xác định và dựa trên nguyên lí toán học. Nếu không phát triển hiểu biết về thủ tục đặt trên nền luận và hiểu, học sinh có thể thực hiện đúng những thủ tục ấy nhưng chỉ xem chúng như một danh sách ``mẹo''. Hệ quả là các em có thể gặp khó khăn khi phải chọn thủ tục thích hợp cho một bài toán cụ thể, hoặc năng lực tưởng như vững vàng trong nhiệm vụ đơn giản lại tan biến khi gặp tình huống phức tạp hơn. Thành thạo thủ tục đích thực đòi hỏi đồng thời việc nắm vững kĩ năng kĩ thuật và phát triển hiểu biết cần thiết để sử dụng chúng một cách thích hợp. Năng lực chiến lược và suy luận thích ứng đều được các thói quen luận trực tiếp đề cập đến.}
}{\NOTE{}}

\TriFigure
  {../figures/fhsm_2009_p12_f1.pdf}
  {../figures/fhsm_2009_p12_f1_vi.pdf}
  {The strands of mathematical proficiency from \emph{Adding It Up} (Kilpatrick, Swafford, and Findell 2001)}
  {Các mạch năng lực toán học theo \emph{Adding It Up}. (Kilpatrick, Swafford, and Findell 2001)}
  {fig:fig_2.1}


\TriPara
  {
    \EN{The development of a ``productive disposition'' (Kilpatrick, Swafford, and Findell 2001) is a high priority of all school mathematics programs. When students achieve this goal, they view mathematics as a reasoning and sense-making enterprise. This goal can be achieved only if students personally engage in mathematical reasoning and sense making as they learn mathematics content.}
  }
  {
    \VI{Việc phát triển ``thái độ tích cực'' đối với toán học (Kilpatrick, Swafford, and Findell 2001) là ưu tiên hàng đầu của mọi chương trình toán học ở trường phổ thông. Khi đạt được mục tiêu này, học sinh xem toán học như một hoạt động dựa trên luận và hiểu. Mục tiêu này chỉ có thể đạt được nếu học sinh trực tiếp tham gia vào luận toán học và hiểu trong quá trình học nội dung toán học.}
  }{\NOTE{}}

\ParallelSection
  {Statistical Reasoning}
  {Luận thống kê}
  {}

  \TriPara
  {
    \EN{Statistical reasoning involves making interpretations based on, and inferences from, data. Statistics furnishes tools for investigating questions by providing strategies for collecting useful data, methods for analyzing the data, and unique perspectives for interpreting the meaning of the data within the context of the problem. The need for statistics arises from “the omnipresence of Conceptual understanding variability” in data (Cobb and Moore 1997, p. 801), and statistical reasoning uses a combination of ideas from both data and chance in seeking to understand that variability. Such notions as distribution, center, spread, association, uncertainty, randomness, sampling, and statistical experiments are foundational concepts underlying the development of statistical reasoning.}
  }
  {
    \VI{Luận thống kê bao gồm việc đưa ra diễn giải dựa trên dữ liệu và suy luận từ dữ liệu. Thống kê cung cấp công cụ để khảo sát câu hỏi bằng những chiến lược thu thập dữ liệu hữu ích, phương pháp phân tích dữ liệu, và góc nhìn đặc thù để diễn giải ý nghĩa của dữ liệu trong bối cảnh bài toán. Nhu cầu đối với thống kê nảy sinh từ ``sự hiện diện khắp nơi của tính biến thiên'' trong dữ liệu (Cobb and Moore 1997, p. 801), và luận thống kê sử dụng kết hợp ý tưởng từ dữ liệu và ngẫu nhiên để tìm cách hiểu tính biến thiên ấy. Những khái niệm như phân bố, vị trí trung tâm, độ phân tán, mối liên hệ, bất định, ngẫu nhiên, lấy mẫu và thí nghiệm thống kê là ý tưởng nền tảng cho sự phát triển của luận thống kê.}
  }{\NOTE{}}


\TriPara
  {
    \EN{Statistics is increasingly recognized as essential for students' success in dealing with the requirements of citizenship, employment, and continuing education (Franklin et al. 2007; College Board 2006, 2007; American Diploma Project 2004). Consequently, the development of statistical reasoning must be a high priority for school mathematics. Garfield (2002) describes a developmental model portraying five stages of statistical reasoning for students, from its lowest level of idiosyncratic reasoning (use of some statistical terms without full understanding) to its highest level of integrated process reasoning (complete understanding of a statistical process). Preparing high school graduates with the ability to make sense of data, and with the capacity to reason with statistics at the integrated process level, requires that students be engaged with meaningful activities involving data and chance throughout prekindergarten through grade 12.}
  }
  {
    \VI{Thống kê ngày càng được công nhận là thiết yếu để học sinh đáp ứng những yêu cầu của đời sống công dân, việc làm và học tập tiếp tục (Franklin et al. 2007; College Board 2006, 2007; American Diploma Project 2004). Do đó, phát triển luận thống kê phải là ưu tiên hàng đầu trong toán học nhà trường. Garfield (2002) mô tả mô hình phát triển gồm năm giai đoạn của luận thống kê ở học sinh, từ mức thấp nhất là luận riêng lẻ theo cảm tính (sử dụng một số thuật ngữ thống kê nhưng chưa hiểu đầy đủ) đến mức cao nhất là luận theo quá trình tích hợp (hiểu đầy đủ một quá trình thống kê). Việc chuẩn bị cho học sinh tốt nghiệp trung học phổ thông khả năng hiểu dữ liệu và năng lực luận bằng thống kê ở mức quá trình tích hợp đòi hỏi các em phải được tham gia vào những hoạt động có ý nghĩa liên quan đến dữ liệu và ngẫu nhiên xuyên suốt từ bậc tiền mẫu giáo đến lớp 12.}
  }{\NOTE{}}

\ParallelSection
  {Mathematical Modeling}
  {Mô hình hóa toán học}
  {}

\TriPara
  {
    \EN{The tools and reasoning processes of mathematics help us understand and operate in the physical and social worlds. Mathematical modeling involves a process of connecting mathematics with a real-world situation; figure 2.2 outlines a cycle often used to organize reasoning in mathematical modeling. A mathematical model is essentially an axiomatization of some sliver of the real world to be able to deal with it mathematically. The connections between mathematics and real-world problems developed in mathematical modeling add value to, and provide incentive and context for, studying mathematical topics.}
  }
  {
    \VI{Công cụ và quá trình luận của toán học giúp chúng ta hiểu và hoạt động trong thế giới vật lí và xã hội. Mô hình hóa toán học bao hàm quá trình kết nối toán học với tình huống trong thế giới thực; Hình 2.2 phác họa chu trình thường được dùng để tổ chức luận trong mô hình hóa toán học. Về bản chất, mô hình toán học là sự tiên đề hóa một lát cắt nào đó của thế giới thực, để có thể xử lí lát cắt ấy bằng toán học. Những kết nối giữa toán học và bài toán trong thế giới thực được phát triển trong mô hình hóa toán học vừa làm cho việc học chủ đề toán học có thêm giá trị, vừa tạo ra động lực và bối cảnh cho việc học ấy.}
  }
  {
    \NOTE{Đoạn này phân biệt rõ giữa \emph{mô hình hóa toán học} và \emph{mô hình toán học}. Mô hình hóa là quá trình nối toán học với tình huống trong thế giới thực và tổ chức luận trong quá trình ấy; mô hình toán học là sản phẩm được tạo ra để xử lí một phần của thế giới thực bằng công cụ toán học. Vì vậy, tiêu đề mục cần giữ là ``Mô hình hóa toán học'', không rút thành ``Mô hình toán học''.}
  }

\TriFigure
  {../figures/fhsm_2009_p13_f1.pdf}
  {../figures/fhsm_2009_p13_f1_vi.pdf}
  {Four-step modeling cycle used to organize reasoning about mathematical modeling}
  {Chu trình mô hình hóa bốn bước dùng để tổ chức luận trong mô hình hóa toán học}
  {fig: fig_2.2}

\TriPara
  {
    \EN{Several additional benefits can be realized from modeling, as can be seen in several examples in section 2. First, mathematical modeling also offers opportunities to make reasoned connections among different mathematical arenas, because in many situations, real-world problems require a combination of mathematical tools. Example 17, ``Clearing the Bridge,'' is a vivid example of how mathematical modeling can cross several mathematical strands. Second, modeling gives students an opportunity to combine mathematical ideas in novel ways. Not only is this ability a life skill, but the ability to use knowledge effectively in new situations has been highlighted in several reports on the workforce and innovation (Partnership for 21st Century Skills 2007; Secretary's Commission on Achieving Necessary Skills 1991; Association for Operations Management [APICS] 2001). Third, working to establish a mathematical model can allow students to become engrossed in the mathematics in ways that promote mathematical reasoning. Fourth, modeling can provide a contextual need to develop mathematical ideas as well as apply them. For example, creating an animation by modeling the movement of images in a plane can serve as motivation to introduce and study matrix theory. Fifth, modeling situations can provide points of access for learners with various backgrounds and skills. Many modeling examples can be continued to very advanced levels. For instance, example 17, ``Clearing the Bridge,'' can be extended into the calculus classroom. Finally, as exemplified by example 14, ``Picture This,'' some modeling contexts can serve as the bases of several lessons-reflecting the genuine fact that persistent, extended efforts are sometimes needed to solve mathematics problems.}
  }
  {
    \VI{Nhiều lợi ích bổ sung có thể thu được từ mô hình hóa, như có thể thấy qua nhiều ví dụ trong Phần 2. Thứ nhất, mô hình hóa toán học cũng tạo cơ hội thiết lập những kết nối có cơ sở luận giữa các mảng toán học khác nhau, bởi trong nhiều tình huống, bài toán thực tế đòi hỏi sự kết hợp nhiều công cụ toán học. Ví dụ 17, ``Chui cầu'', là minh họa sinh động cho thấy mô hình hóa toán học có thể đi qua nhiều mạch toán học. Thứ hai, mô hình hóa cho học sinh cơ hội kết hợp ý tưởng toán học theo những cách mới mẻ. Năng lực này không chỉ là kĩ năng sống; khả năng sử dụng tri thức hiệu quả trong tình huống mới còn được nhấn mạnh trong nhiều báo cáo về lực lượng lao động và đổi mới (Partnership for 21st Century Skills 2007; Secretary's Commission on Achieving Necessary Skills 1991; Association for Operations Management [APICS] 2001). Thứ ba, nỗ lực xây dựng một mô hình toán học có thể khiến học sinh cuốn sâu vào toán học theo cách thúc đẩy luận toán học. Thứ tư, mô hình hóa có thể tạo ra nhu cầu theo ngữ cảnh để phát triển ý tưởng toán học cũng như áp dụng chúng. Ví dụ, việc tạo hoạt hình bằng cách mô hình hóa chuyển động của hình ảnh trong mặt phẳng có thể trở thành động lực để giới thiệu và nghiên cứu lí thuyết ma trận. Thứ năm, tình huống mô hình hóa có thể tạo điểm tiếp cận cho người học có nền tảng và kĩ năng khác nhau. Nhiều ví dụ mô hình hóa có thể được tiếp tục phát triển đến mức rất cao. Chẳng hạn, Ví dụ 17, ``Chui cầu'', có thể được mở rộng vào lớp học giải tích. Cuối cùng, như được minh họa qua Ví dụ 14, ``Đứng đâu để chụp?'', một số bối cảnh mô hình hóa có thể làm nền cho nhiều bài học, phản ánh thực tế rằng đôi khi cần nỗ lực bền bỉ và kiên trì để giải quyết bài toán.}
  }{\NOTE{}}

\ParallelSection
  {Technology to Support Reasoning and Sense Making}
  {Công nghệ hỗ trợ luận và hiểu}
  {}


\TriPara
  {
    \EN{Technology is an integral part of society, the workplace, and even many areas of modern mathematical research itself (Bailey and Borwein 2005); the mathematics classroom needs to reflect that reality. Technology can be used to advance the goals of reasoning and sense making in the high school mathematics classroom. Technological tools can be particularly useful in looking for patterns and relationships and in forming conjectures---see examples 9, ``More Than Meets the Eye''; 14, ``Picture This''; and 17, ``Clearing the Bridge.'' Technology can relieve students of burden some computations, giving them the freedom, and the need, to think strategically, as in example 9, ``More Than Meets the Eye.'' Using technology to display multiple representations of the same problem can aid in making connections, as in example 11, ``Take As Directed.'' When technology allows multiple representations to be linked dynamically, it can provide new opportunities for students to take mathematically meaningful actions and immediately see mathematically meaningful consequences---fertile ground for sense-making and reasoning activities. This dynamic linking is evident in example 13, ``Tidal Waves.'' Technology can also be useful in generalizing a solution, as in example 22, ``Meaningful Words, Part B.''}
  }
  {
    \VI{Công nghệ là một phần không thể tách rời của xã hội, nơi làm việc, và thậm chí nhiều lĩnh vực nghiên cứu toán học hiện đại (Bailey and Borwein 2005); lớp học toán cần phản ánh thực tế đó. Công nghệ có thể được sử dụng để thúc đẩy mục tiêu luận và hiểu trong lớp học toán trung học phổ thông. Công cụ công nghệ có thể đặc biệt hữu ích khi tìm kiếm mẫu hình và quan hệ, cũng như khi hình thành phỏng đoán---xem Ví dụ 9, ``Chưa thấy hết quy luật''; Ví dụ 14, ``Đứng đâu để chụp?''; và Ví dụ 17, ``Chui cầu''. Công nghệ có thể giúp học sinh giảm bớt gánh nặng tính toán, cho các em sự tự do, đồng thời tạo ra nhu cầu tư duy chiến lược, như trong Ví dụ 9, ``Chưa thấy hết quy luật''. Việc sử dụng công nghệ để hiển thị nhiều biểu diễn của cùng một bài toán có thể hỗ trợ thiết lập kết nối, như trong Ví dụ 11, ``Uống theo chỉ dẫn''. Khi công nghệ cho phép nhiều biểu diễn được liên kết động, nó có thể tạo ra cơ hội mới để học sinh thực hiện hành động có ý nghĩa toán học và ngay lập tức thấy được hệ quả có ý nghĩa toán học---mảnh đất màu mỡ cho hoạt động hiểu và luận. Sự liên kết động này được thể hiện trong Ví dụ 13, ``Thủy triều''. Công nghệ cũng có thể hữu ích trong việc khái quát hóa lời giải, như trong Ví dụ 22, ``Từ có nghĩa dễ nhớ hơn! (B)''.}
  }
  {
    \NOTE{Điểm then chốt là liên kết động giữa các biểu diễn: khi biểu diễn toán học được kết nối và phản hồi tức thời với nhau, học sinh có thể thực hiện thao tác có ý nghĩa toán học và quan sát ngay hệ quả tương ứng. Chính sự tương tác này tạo ra môi trường thuận lợi để hiểu và luận diễn ra tự nhiên.}
  }

\TriPara
  {
    \EN{The incorporation of technology in the classroom should not overshadow the development of the procedural proficiency needed by students to support continued mathematical growth. It should be used as a tool that leads to a deeper understanding of mathematical concepts. Students can be challenged to take responsibility for deciding which tool might be useful in a given situation when they are allowed to choose from a menu of mathematical tools that includes technology. Students who have regular opportunities to discuss and reflect on how a  technological tool is used effectively will be less apt to use technology as a crutch.}
  }
  {
    \VI{Việc đưa công nghệ vào lớp học không nên làm lu mờ quá trình phát triển thành thạo thủ tục mà học sinh cần để tiếp tục tăng trưởng toán học. Công nghệ cần được sử dụng như công cụ dẫn đến hiểu biết sâu hơn về khái niệm toán học. Khi được phép lựa chọn từ danh mục công cụ toán học bao gồm cả công nghệ, học sinh có thể được thử thách để tự chịu trách nhiệm quyết định công cụ nào hữu ích trong tình huống đang xét. Những học sinh thường xuyên có cơ hội thảo luận và phản tư về cách sử dụng hiệu quả công cụ công nghệ sẽ ít có xu hướng dùng công nghệ như chỗ dựa thay cho suy nghĩ.}
  }{\NOTE{}}

\ParallelSection
  {Conclusion}
  {Kết luận}
  {}

\TriPara
  {
    \EN{Reasoning and sense making must become a part of the fabric of the high school mathematics classroom. Not only are they important goals themselves, but they are the foundation for true mathematical competence. Incorporating isolated experiences with reasoning and sense making will not suffice. Teachers must consistently support and encourage students' progress toward more sophisticated levels of reasoning.}
  }
  {
    \VI{Luận và hiểu phải trở thành một phần trong kết cấu của lớp học toán trung học phổ thông. Chúng không chỉ là những mục tiêu quan trọng tự thân, mà còn là nền tảng của năng lực toán học đích thực. Việc chỉ lồng ghép những trải nghiệm rời rạc liên quan đến luận và hiểu là không đủ. Giáo viên cần nhất quán hỗ trợ và khuyến khích sự tiến triển của học sinh hướng tới những mức độ luận ngày càng tinh vi hơn.}
  }{\NOTE{}}

\CloseGrid